\addtocontents{toc}{\protect\Needspace{3\baselineskip}}
\originalchapter{数值宇宙学项目：从粒子快照到晕目录}
8.902 Fall 2023 还公开了一份五页数值项目说明，任务是运行 Gadget-4 并分析 $z=0$ 快照。公开文件给出了项目目标、算法参数和预期输出；课堂 Canvas 上的预打包配置、初始条件生成设置、读入脚本与实际快照并未作为公开资料取得。本章完整交代公开题目所要求的计算、实现检查点与可推导的结果形式。没有快照便不能报告“第三大晕”的数值、拟合浓度或真实质量函数，也不能把示意数据冒充该次作业的运行结果。

\originalsection{题目的数据路径和物理目标}
项目要求用 Gadget-4 的构型文件、参数文件、快照时刻表完成一次宇宙学 $N$ 体运行，再用 Python 后处理。课程预设的运行示例在同步输出中出现 $884736$ 个引力粒子；盒边长为 $40\,\mathrm{Mpc}$，粒子质量相同。它要求读入三个时间快照，并取最后的 \texttt{snapshot\_002.hdf5}，该文件对应 $z=0$。原项目说明说初始条件在运行时由配置产生，因此单独下载一个 Gadget-4 发布版并不足以再现这份课程配置。编译需要 C、\texttt{make}、FFTW、GSL、HDF5、MPI；实际参数单位要读随课程配置的 \texttt{System of Units}，不可把 HDF5 中所有无说明字段直接当作物理单位。

\src{8.902 Fall 2023, Hands-on Numerical Project, pp.1--4；课堂 Canvas 数据依赖未取得。}
\begin{exercise}
对 $z=0$ 粒子快照：(a) 用 friends-of-friends (FoF) 算法构建晕目录，链接长度取平均粒子间距的 $0.2$ 倍；(b) 因两个最重晕正在合并，选择第三大晕，绘制其 $30\,\mathrm{kpc}$ 至 $1\,\mathrm{Mpc}$ 径向密度并与 NFW 模型比较，估计 $R_{200,\mathrm{crit}}$ 与浓度；(c) 仅对至少含 20 个粒子的晕制作质量直方图。课程要求三人组在12月7日做约15分钟展示，没有书面提交件。
\end{exercise}
\begin{solution}
设盒边长 $L$、总粒子数 $N$，平均粒子间距是 $\bar d=(L^3/N)^{1/3}$，FoF 阈值 $\ell=0.2\bar d$；若实际快照恰有 $N=884736$ 且 $L=40\,\mathrm{Mpc}$，则 $\bar d\simeq0.417\,\mathrm{Mpc}$、$\ell\simeq0.0834\,\mathrm{Mpc}$。这里 $884736$ 来自运行说明的示例输出，仍须用所取得的真实快照头部核验，不能把这个条件值说成已读文件。每对周期边界粒子用最小镜像距离
\[
 \Delta x_j\leftarrow\Delta x_j-L\,\operatorname{round}(\Delta x_j/L),\qquad
 |\boldsymbol{\Delta x}|<\ell
\]
建立无向边；FoF 晕是图的连通分量，不是距某个固定中心都小于 $\ell$ 的球。可用空间网格或邻近搜索枚举可能的边，再用并查集求分量。晕质量 $M_i=N_i m_p$；按 $M_i$ 递减取第三个，需说明并列、粒子类型与质量单位的规则。大 FoF 组可能用细桥连接，给密度中心和 $R_{200}$ 时应重新定义中心，不将 FoF 外缘当作 $R_{200}$。

以该晕中心为零点，周期展开成员相对位移。半径壳 $[r_j,r_{j+1})$ 的直接估计为
\[
 \hat\rho_j=\frac{m_pN_j}{(4\pi/3)(r_{j+1}^3-r_j^3)} .
\]
内区还要核对力软化和粒子数，外区不能把盒平均背景错误计入单个晕；课程规定的比较区间是 $0.03$ 至 $1\,\mathrm{Mpc}$。NFW 形式为
\[
 \rho(r)=\frac{\rho_s}{(r/r_s)(1+r/r_s)^2},\quad
 M(<r)=4\pi\rho_s r_s^3\left[\ln(1+x)-\frac{x}{1+x}\right],\quad x=r/r_s .
\]
定义 $\bar\rho(<R_{200,\mathrm{crit}})=200\rho_{\mathrm{crit}}(z=0)$，其中 $\rho_{\mathrm{crit}}=3H_0^2/(8\pi G)$；再给 $c_{200}=R_{200,\mathrm{crit}}/r_s$。若用对数密度拟合，须记录壳宽、误差权重、排除半径和同一单位系统。成员数至少 20 的晕进入 $M_i=N_im_p$ 的分箱；该\emph{直方图}依箱宽而变，若称质量函数则应另除盒体积及相应的 $\Delta M$ 或 $\Delta\ln M$。

完整运行的必要交叉检查是：头部粒子数与质量和、周期链接的不依赖盒切面性、FoF 分量计数不重复、径向壳质量之和与选择的成员质量相符、$M(<r)$ 单调及 $R_{200}$ 定义自洽。以上给出算法和可验证方程；没有公开的该次快照，故无该晕的实际曲线、$r_s$、$R_{200}$、$c_{200}$ 或晕质量数值。
\end{solution}
